\documentclass[10pt,twocolumn,a4paper]{article}
\usepackage[utf8]{inputenc}
\usepackage[T1]{fontenc}
\usepackage{lmodern}
\usepackage{microtype}
\usepackage[margin=2.2cm]{geometry}
\usepackage{amsmath,amssymb,amsthm,mathtools}
\usepackage{bm}
\usepackage{physics}
\usepackage{graphicx}
\usepackage{booktabs}
\usepackage{array}
\usepackage{multirow}
\usepackage{tabularx}
\usepackage{enumitem}
\usepackage{float}
\usepackage{caption}
\usepackage{subcaption}
\usepackage{xcolor}
\usepackage{siunitx}
\usepackage[colorlinks=true,linkcolor=blue!60!black,
            citecolor=blue!60!black,urlcolor=blue!60!black]{hyperref}
\usepackage{cleveref}

\captionsetup{font=small,labelfont=bf,textfont=it}

\newtheorem{theorem}{Theorem}[section]
\newtheorem{lemma}[theorem]{Lemma}
\newtheorem{proposition}[theorem]{Proposition}
\newtheorem{corollary}[theorem]{Corollary}
\theoremstyle{definition}
\newtheorem{definition}[theorem]{Definition}
\newtheorem{remark}[theorem]{Remark}

\newcommand{\Hb}{\mathrm{Hb}}
\newcommand{\HbO}{\mathrm{HbO_2}}
\newcommand{\SpO}{\mathrm{SpO_2}}
\newcommand{\Tskin}{T_{\mathrm{skin}}}
\newcommand{\HRobs}{\mathrm{HR}_{\mathrm{obs}}}
\newcommand{\HRint}{\mathrm{HR}_{\mathrm{int}}}
\newcommand{\dHRm}{\Delta\mathrm{HR}_{\mathrm{met}}}
\newcommand{\dHRo}{\Delta\mathrm{HR}_{\mathrm{O_2}}}
\newcommand{\dHRa}{\Delta\mathrm{HR}_{\mathrm{auto}}}
\newcommand{\Qten}{Q_{10}}
\newcommand{\vaso}{\eta}
\newcommand{\melan}{m}
\newcommand{\Hbconc}{[\Hb]_{\mathrm{tot}}}
\newcommand{\Soxy}{S_{\mathrm{O_2}}}

\title{\textbf{Layered Optical Inversion for Photoplethysmography:\\
Recovering Vasodilation, Skin Temperature, and a Metabolic--Autonomic\\
Decomposition of Heart Rate from RGB Camera Pixels}\\[0.3em]}

\author{
Kundai Farai Sachikonye\\[0.3em]
AIMe Registry for Artificial Intelligence\\
Experimental Bioinformatics, Maximus-von-Imhof-Forum 3\\
85354 Freising, Germany\\
\texttt{kundai.sachikonye@bitspark.com}
}

\date{\today}

\begin{document}

\twocolumn[
\begin{@twocolumnfalse}
\maketitle
\begin{abstract}
\noindent
We present a forward-model-based inversion method for camera-based
photoplethysmography (rPPG) that recovers four optical state
parameters---melanin index, total hemoglobin concentration, oxygenation,
and an effective vasodilation factor---from the time-averaged RGB pixel
values inside a face region of interest, and that converts the
vasodilation factor into a skin-temperature estimate and a metabolic
decomposition of observed heart rate. Conventional rPPG extracts
heart rate from temporal modulation of skin reflectance (typically the
green channel) without modelling the underlying optical pathway, and
therefore cannot separate cardiac drive from thermal drive, hypoxic
drive, or skin-pigmentation-driven baseline shifts. We instead model
the forehead as a layered absorber--scatterer composed of an
air--surface interface, an epidermis with melanin absorption, and a
dermis with hemoglobin in a capillary network whose effective optical
path length is the vasodilation factor. Forward propagation through
the stack via the Beer--Lambert law gives predicted RGB reflectances as
a function of the four state parameters. The inverse problem admits
algebraic solution via channel separation: the blue band identifies
melanin, red identifies hemoglobin given melanin, and green identifies
oxygenation given both. Vasodilation is read from the AC modulation
amplitude of the pulsatile component, normalised against a 30~s
rolling baseline. We map vasodilation to skin temperature through a
linearised thermoregulation relation, then use the Arrhenius
$\Qten$~law to decompose the observed heart rate into intrinsic,
metabolic, hypoxic, and autonomic components, isolating the genuine
neural drive from purely thermal contributions. We validate the method
through six simulation experiments. Self-consistency demonstrates
that the iterative algebraic inverse recovers melanin to median
$0.5\%$ relative error and total hemoglobin concentration to
median $10\%$ in the noise-free limit; oxygenation recovery from
RGB alone has an inherent ceiling around $40\%$ relative error owing
to the small green-channel spectral difference between oxy- and
deoxyhemoglobin, and is best treated as a relative-trend signal
rather than an absolute reading. Noise-sensitivity analysis confirms
that melanin and Hb degrade gracefully with measurement noise while
oxygenation error is dominated by structural model approximations.
The vasodilation--temperature mapping responds physiologically across
the range. The PCHR decomposition correctly attributes synthetic
fever, hypoxia, and exercise scenarios to the matching
$\Delta\mathrm{HR}$ component. Under bandpass-rejectable ambient
light drift the layered method's heart-rate variance is comparable
to a green-channel-only baseline; the layered method's distinct
value is the additional state recovery (vasodilation, skin
temperature, total hemoglobin) that amplitude-only rPPG does not
produce. The full pipeline runs in real time in a WebGPU-enabled
browser.

\vspace{4pt}
\noindent\textbf{Keywords:} remote photoplethysmography, rPPG, skin optics,
Beer--Lambert, hemoglobin, vasodilation, skin temperature, heart rate
decomposition, $\Qten$~metabolic effect, multi-channel inversion,
camera-based vital sensing.
\end{abstract}
\vspace{12pt}
\end{@twocolumnfalse}
]

% =====================================================================
\section{Introduction}
% =====================================================================

Camera-based photoplethysmography (rPPG) extracts cardiac information
from the temporal modulation of skin reflectance recorded by a
consumer camera~\cite{verkruysse2008,allen2007,mcduff2015}. The
signal arises because each cardiac cycle alters the volume of blood in
dermal microvasculature, modulating the optical absorption of
hemoglobin and producing a small periodic change in reflected light
intensity. Standard methods estimate heart rate from this
modulation by bandpass filtering one or more colour channels (most
commonly green, where hemoglobin absorption is strongest in the
visible band), or by linearly combining channels via algorithms such
as POS~\cite{wang2017posrPPG} or
CHROM~\cite{dehaan2013chrom}.

Conventional rPPG is \emph{pixel-blind} in the sense that the
optical pathway from skin to pixel is not explicitly modelled. The
extracted signal is treated as a one-dimensional time series, with no
attempt to separate the contributions of pigmentation, blood volume,
oxygenation, and ambient illumination. As a result several
biologically meaningful quantities that are in principle present in
the measurement are not recoverable: skin temperature, oxygenation,
and the proportion of heart-rate variability attributable to thermal
versus neural drive. These quantities are not exotic; they are the
inputs to standard cardiovascular physiology
models~\cite{guyton2006,levick2009} and would normally be obtained by
separate sensors (thermistor, pulse oximeter, ECG with heart-rate
variability analysis). Recovering them from the same camera that
already produces the heart-rate signal would substantially extend the
diagnostic utility of consumer rPPG.

In this paper we develop a forward optical model of forehead skin and
its inverse, and use the resulting state estimates to decompose
observed heart rate into intrinsic, metabolic, hypoxic, and autonomic
components. The forward model is a layered Beer--Lambert
absorber--scatterer with four state parameters: melanin index
($\melan$), total hemoglobin concentration ($\Hbconc$), arterial
oxygenation fraction ($\Soxy$), and an effective vasodilation factor
($\vaso$). The inverse problem admits closed-form algebraic solution
via channel separation: the blue band is dominated by melanin, red by
hemoglobin given fixed melanin, green by oxygenation given fixed
melanin and hemoglobin, and the AC modulation amplitude of the
pulsatile component supplies the vasodilation factor independently of
the DC channels.

We then map vasodilation to skin temperature through a linearised
thermoregulatory relation calibrated against published forehead
temperature ranges, and use the Arrhenius $\Qten$~law to decompose
heart rate into:
\begin{equation}
    \HRobs \;=\; \HRint + \dHRm + \dHRo + \dHRa,
\end{equation}
where $\HRm$ is the thermal drive (metabolic-rate-mediated), $\dHRo$
the hypoxic compensation, and $\dHRa$ the autonomic
residual---the part of heart rate that genuinely reflects neural
drive. Existing rPPG implementations conflate these terms; the present
method separates them.

\paragraph{Contributions.} Specifically:
\begin{enumerate}[leftmargin=*,nosep]
\item A closed-form layered Beer--Lambert forward model for forehead
  skin reflectance whose state space is $(\melan, \Hbconc, \Soxy,
  \vaso)$ and whose output is RGB camera reflectance with
  channel-band-averaged absorption coefficients calibrated from
  standard tissue optical data.
\item An algebraic inverse solution by channel separation that
  recovers all four state parameters in $\sim$100~floating-point
  operations per face ROI per frame, with provable identifiability
  under physiological parameter ranges.
\item A vasodilation-to-temperature mapping derived from
  thermoregulatory linearisation, and a Q$_{10}$-based decomposition
  of heart rate into intrinsic, metabolic, hypoxic, and autonomic
  components.
\item A real-time WebGPU implementation that produces
  $1\,\mathrm{Hz}$ state estimates from a 720p webcam.
\item Six simulation experiments validating self-consistency,
  noise sensitivity, identifiability, temperature mapping, scenario
  attribution, and lighting robustness against a green-only baseline.
\end{enumerate}

The paper is organised as follows.
\Cref{sec:background} reviews the relevant background in skin optics
and rPPG.
\Cref{sec:forward} derives the forward model.
\Cref{sec:inverse} states and solves the inverse problem.
\Cref{sec:cardiac} converts the optical state into a metabolic
decomposition of heart rate.
\Cref{sec:implementation} describes the real-time implementation.
\Cref{sec:methods} states the experimental methods used for
validation.
\Cref{sec:results} reports the results.
\Cref{sec:discussion} discusses limitations and future directions, and
\cref{sec:conclusion} concludes.

% =====================================================================
\section{Background}\label{sec:background}
% =====================================================================

\subsection{Optical Structure of Forehead Skin}\label{ssec:skin}

The forehead is one of the most thoroughly characterised body sites
for rPPG because its skin is relatively flat, sparsely
hair-covered, and highly vascularised in the papillary
dermis~\cite{anderson1981skin,jacques2013}. From the camera's
viewpoint, the forehead presents an approximately laterally uniform
stack of optical layers:

\begin{itemize}[leftmargin=*,nosep]
\item \textbf{Stratum corneum} ($10$--$40\,\mu$m): outermost dead
  cell layer, predominantly keratin and lipid; refractive index
  $n_{\mathrm{sc}}\approx 1.55$~\cite{tuchin2007}; minimal absorption
  in the visible band; surface microbial film modifies effective
  refractive index slightly~\cite{prescott2017}.
\item \textbf{Epidermis} ($50$--$150\,\mu$m): living cell layer
  containing melanin in melanosomes; melanin absorbs broadly across
  the visible spectrum with monotonically decreasing absorption from
  blue to red; scattering is moderate~\cite{jacques1998melanin}.
\item \textbf{Dermis} ($1$--$3\,\mathrm{mm}$): connective tissue
  containing the capillary network responsible for the rPPG signal;
  hemoglobin is the dominant chromophore, with strong absorption peaks
  for both oxyhemoglobin (\HbO, $\lambda_{\max}\approx
  542,576\,\mathrm{nm}$) and deoxyhemoglobin (\Hb,
  $\lambda_{\max}\approx
  555\,\mathrm{nm}$)~\cite{prahlomlc,zijlstra2000}.
\item \textbf{Subcutaneous adipose} ($\geq 1\,\mathrm{mm}$): mostly
  lipid; high diffuse backscatter contributes to the apparent
  reflectance baseline.
\end{itemize}

The cardiac modulation is contributed almost entirely by changes in
dermal blood volume during the cardiac cycle. Each ventricular
ejection drives a small increase in capillary volume; the consequent
increase in path length through hemoglobin reduces reflectance. The
modulation depth is on the order of $0.1$--$1\%$ at typical
illumination~\cite{verkruysse2008}, which is at the limit of
useful camera bit-depth and is therefore highly sensitive to
illumination noise.

\subsection{Hemoglobin Spectroscopy}

The molar extinction coefficients of \HbO\ and \Hb\ are well
characterised from
$250$--$1000\,\mathrm{nm}$~\cite{prahlomlc,kim2005}. For the present
work we require band-averaged values over the camera's
red, green, and blue channel responses (centred near $620$, $540$,
$465\,\mathrm{nm}$ respectively, with bandwidths around $80\,\mathrm{nm}$).
Integrated against typical CMOS quantum-efficiency curves, the
band-averaged molar extinction values are approximately
$\varepsilon_{\Hb,R}\approx 770\,\mathrm{M^{-1}cm^{-1}}$,
$\varepsilon_{\Hb,G}\approx 5300\,\mathrm{M^{-1}cm^{-1}}$,
$\varepsilon_{\Hb,B}\approx 8900\,\mathrm{M^{-1}cm^{-1}}$, with
analogous values for \HbO~\cite{prahlomlc}. Total
absorption depends on hemoglobin concentration, oxygenation, and
optical path length through dermal blood. The latter is what
vasodilation modulates.

\subsection{Conventional rPPG}

The first systematic camera-based PPG demonstration is due to
Verkruysse, Svaasand and Nelson~\cite{verkruysse2008}, who showed that
ordinary webcam recordings carry a heart-rate signal recoverable by
bandpass filtering the green channel of a face ROI. Subsequent work
introduced multi-channel methods that improve robustness to motion
and lighting: independent component analysis~\cite{poh2010},
chrominance-based projection (CHROM)~\cite{dehaan2013chrom}, and
plane-orthogonal-to-skin (POS)~\cite{wang2017posrPPG}. These methods
share the property that they extract heart rate from the temporal
amplitude of skin reflectance modulation but do not invert the
underlying biophysical model. Several extensions estimate respiration
rate~\cite{tarassenko2014} or oxygenation~\cite{verkruysse2009} as
secondary signals, but the field remains dominated by purely
algorithmic time-series approaches.

\subsection{Heart-Rate Components in Physiology}

Standard cardiovascular physiology
texts~\cite{guyton2006,berne2010,levick2009} decompose observed heart
rate into an intrinsic pacemaker rate (typically $90$--$110\,\mathrm{bpm}$
in the denervated sino-atrial node) modified by metabolic
state, hypoxic compensation, and autonomic balance. Body-temperature
elevation increases heart rate via the Arrhenius $\Qten$~effect on
sino-atrial node automaticity, with $\Qten\approx 2.3$ for cardiac
pacemaker activity~\cite{stevens1989q10,delgado2008}. Hypoxic compensation acts via
chemoreceptor reflexes that increase sympathetic outflow when arterial
oxygenation drops~\cite{guyenet2014}. The autonomic balance reflects
the difference between sympathetic activation and parasympathetic
withdrawal and is the component of interest in stress, attention, and
clinical applications. These three drives are conflated in
amplitude-based rPPG: a warm and resting subject and a normothermic
and stressed subject can produce identical raw heart rates. Separating
them requires either independent skin-temperature and oxygenation
sensors, or a model that derives them from the same camera data.

% =====================================================================
\section{Layered Optical Forward Model}\label{sec:forward}
% =====================================================================

\subsection{Geometry and Notation}

We treat the forehead as a slab geometry with light incident
approximately normally from the camera direction, propagating through
four layers indexed $j\in\{0,1,2,3\}$, and back-diffusing toward the
camera (\cref{fig:layered-model}). Layer~0 is the air--skin interface;
layer~1 is the stratum corneum; layer~2 the epidermis; layer~3 the
dermis with capillary blood; the deeper subcutaneous fat is treated as
a Lambertian back-reflector with effective albedo $D$. Let
$\rho_{j}(\lambda)$ denote the transmittance of layer~$j$ at wavelength
$\lambda$. The observed reflectance per wavelength is
\begin{equation}
    R(\lambda)
    \;=\; \rho_{\mathrm{spec}}(\lambda)
    + \bigl[1 - \rho_{\mathrm{spec}}(\lambda)\bigr]
      \prod_{j=1}^{3}\rho_{j}^{2}(\lambda)\, D,
    \label{eq:R-of-lambda}
\end{equation}
where the squared transmittance accounts for the down-and-back photon
path, $\rho_{\mathrm{spec}}$ is the surface specular reflectance set
by Fresnel and surface microstructure, and $D$ is the diffuse return
of the subcutaneous boundary.



\subsection{Layer Transmittances}

\paragraph{Stratum corneum (layer~1).} For the visible band, absorption
in layer~1 is negligible relative to the other layers; we set
$\rho_{1}(\lambda)\approx 1$. Surface microbial colonisation and
sebum layer thickness alter the specular term slightly; we collect
these effects into a constant correction
$\delta_{\mu}\in[-0.04,+0.04]$ on $\rho_{\mathrm{spec}}$ that is held
fixed unless multi-day reflectance data are available for fit.

\paragraph{Epidermis (layer~2).} Layer~2 transmittance is dominated by
melanin, with absorption coefficient
$\mu_{a,\mathrm{mel}}(\lambda)$ approximately following a
$\lambda^{-3.4}$ power law in the visible
band~\cite{jacques1998melanin}. We parametrise the melanin content by
a dimensionless index $\melan\in[0,1]$ such that, for layer thickness
$L_{2}$ and channel-band-averaged absorption coefficient
$\bar{\mu}_{\mathrm{mel},c}$ in channel $c\in\{R,G,B\}$,
\begin{equation}
    \rho_{2,c}
    \;=\; \exp\bigl(-\bar{\mu}_{\mathrm{mel},c}\, \melan\, L_{2}\bigr).
    \label{eq:rho2}
\end{equation}
Numerical values of $\bar{\mu}_{\mathrm{mel},c} L_{2}$ are computed
from the band-averaged Jacques curve and tabulated in
\cref{tab:abs-coeffs}.

\paragraph{Dermis (layer~3).} Layer~3 transmittance is dominated by
hemoglobin in the capillary network. Letting $\bar{\mu}_{\Hb,c}$ and
$\bar{\mu}_{\HbO,c}$ denote the band-averaged extinction coefficients
of deoxy- and oxy-hemoglobin (per unit total hemoglobin concentration
in $\mathrm{g/dL}$), and writing $\Soxy$ for the oxygenation fraction,
the dermal absorption coefficient at channel $c$ is
\begin{equation}
    \mu_{a,3,c}
    \;=\; \Hbconc\,\bigl[(1-\Soxy)\bar{\mu}_{\Hb,c} + \Soxy\bar{\mu}_{\HbO,c}\bigr].
    \label{eq:mu-derm}
\end{equation}
The corresponding transmittance is
\begin{equation}
    \rho_{3,c}
    \;=\; \exp\bigl(-\mu_{a,3,c}\, \vaso\, L_{3,0}\bigr),
    \label{eq:rho3}
\end{equation}
where $L_{3,0}$ is the resting effective dermal optical path length
and $\vaso\in[0.6,1.6]$ is the dimensionless vasodilation factor that
modulates effective path length: vasoconstriction reduces capillary
volume and effective path; vasodilation increases both. We define
the resting state $\vaso=1$ to coincide with a forehead skin
temperature of $\Tskin=33\,^{\circ}\mathrm{C}$ at quiet
sitting~\cite{webb1992,kelechava2014}.

\subsection{Channel-Band-Averaged Forward Model}

Substituting~\cref{eq:rho2,eq:rho3} into~\cref{eq:R-of-lambda} and
collecting per-channel constants into the parameters of
\cref{tab:abs-coeffs} yields the channel-band-averaged forward map
\begin{equation}
    \boxed{\;
    R_{c}(\theta)
    \;=\; \rho_{\mathrm{spec}}
    + (1 - \rho_{\mathrm{spec}})\,
      \rho_{2,c}^{2}(\melan)\,
      \rho_{3,c}^{2}(\Hbconc,\Soxy,\vaso)\,
      D,
    \;}
    \label{eq:forward}
\end{equation}
for $c\in\{R,G,B\}$, where $\theta=(\melan,\Hbconc,\Soxy,\vaso)$ is
the four-dimensional state. Equation~\eqref{eq:forward} is the forward
model used throughout this paper.

\begin{table}[H]
\centering
\caption{Channel-band-averaged effective absorption coefficients used
in~\cref{eq:rho2,eq:rho3}. Values absorb the layer thicknesses
$L_{2}$, $L_{3,0}$ to give dimensionless products. Sources:
melanin~\cite{jacques1998melanin}, hemoglobin~\cite{prahlomlc}.}
\label{tab:abs-coeffs}
\small
\begin{tabular}{lccc}
\toprule
& Red ($\sim$620\,nm) & Green ($\sim$540\,nm) & Blue ($\sim$465\,nm) \\
\midrule
$\bar{\mu}_{\mathrm{mel},c} L_{2}$ & 0.10 & 0.36 & 0.84 \\
$\bar{\mu}_{\Hb,c}\, L_{3,0}$ (per g/dL) & 0.012 & 0.112 & 0.084 \\
$\bar{\mu}_{\HbO,c}\, L_{3,0}$ (per g/dL) & 0.012 & 0.148 & 0.084 \\
\bottomrule
\end{tabular}
\end{table}

\subsection{Surface Specular and Diffuse Constants}

We use $\rho_{\mathrm{spec}}=0.04$, consistent with Fresnel reflection
at the air--skin interface for moderate angles, and $D=0.55$,
consistent with diffuse return from subcutaneous fat as reported in
diffuse-reflectance spectroscopy~\cite{tuchin2007,bosschaart2014}.
These are not free parameters; they are fixed for all subjects unless
specifically calibrated.

\subsection{Vasodilation--Temperature Linearisation}

Skin temperature reflects the thermoregulatory state of the
cutaneous microcirculation~\cite{rowell1977,charkoudian2010}. Sympathetic
withdrawal at the cutaneous adrenergic terminals dilates the dermal
arterioles, increases capillary volume, and elevates skin temperature
by routing more warm core blood to the surface. The relationship is
nonlinear in detail but locally linear over the physiological range.
We use
\begin{equation}
    \Tskin = T_{\mathrm{rest}} + \kappa_{T}(\vaso - 1),
    \qquad \kappa_{T} = 4\,\mathrm{K},
    \label{eq:vaso-T}
\end{equation}
with $T_{\mathrm{rest}}=33\,^{\circ}\mathrm{C}$ for resting forehead.
The slope $\kappa_{T}$ is set so that the physiological vasodilation
range $\vaso\in[0.6,1.6]$ maps to
$\Tskin\in[27,37]\,^{\circ}\mathrm{C}$, approximately the
maximum range observed during cold-pressor and exercise
challenges~\cite{webb1992,charkoudian2010}.

\begin{figure*}[!htbp]
    \centering
    \includegraphics[width=\textwidth]{figures/panel4_vaso_temperature.png}
    \caption{\textbf{Vasodilation--Temperature Mapping and the Joint $(\eta, \mathrm{HR}_{\mathrm{int}}, \dHRm)$ Surface --- Metabolic-Drive Geometry.}
    The vasodilation factor recovered from the AC channel is converted to skin temperature via the linearisation in \cref{eq:vaso-T}; the resulting temperature drives the metabolic component of the PCHR decomposition (\cref{eq:dHR-met}).
    (\textbf{A})~Linear mapping $\Tskin = 33 + 4(\eta - 1)\,^{\circ}\mathrm{C}$ across the physiological vasodilation range $[0.6, 1.6]$. The black dot at $(\eta, \Tskin) = (1.0, 33\,^{\circ}\mathrm{C})$ marks the resting reference anchor; the slope $\kappa_{T}=4\,\mathrm{K}$ is calibrated so the full vasodilation range maps onto the empirical thermoregulatory range $[27, 37]\,^{\circ}\mathrm{C}$.
    (\textbf{B})~Distribution of $\Tskin$ values induced by uniformly distributed $\eta \in [0.6, 1.6]$ ($n=5000$ samples). The histogram is approximately uniform across $\sim 27$--$37\,^{\circ}\mathrm{C}$, the support over which the metabolic decomposition is exercised in \cref{fig:pchr-scenarios}.
    (\textbf{C})~Three-dimensional surface of $\dHRm$ as a function of $\eta$ and intrinsic heart rate $\mathrm{HR}_{\mathrm{int}}$, shaded by the metabolic-drive magnitude on a plasma colour map. The surface is bilinear in $(\eta - 1)$ and $\mathrm{HR}_{\mathrm{int}}$ as expected from the product structure of \cref{eq:dHR-met}; the maximum drive is reached at the high-$\eta$ / high-$\mathrm{HR}_{\mathrm{int}}$ corner. Negative drive at $\eta < 1$ reflects vasoconstriction below the thermal-reference state.
    (\textbf{D})~$\dHRm$ versus $\Tskin$ at four intrinsic heart rates ($50,60,70,80\,\mathrm{bpm}$). All curves cross zero at $T_{\mathrm{rest}}=33\,^{\circ}\mathrm{C}$ with slope $\alpha_{T}\,\mathrm{HR}_{\mathrm{int}}$. A subject with $\mathrm{HR}_{\mathrm{int}}=60$ at $\Tskin=35\,^{\circ}\mathrm{C}$ contributes $\sim 9.6\,\mathrm{bpm}$ of metabolic drive --- the typical thermal acceleration during a warm afternoon.}
    \label{fig:vaso-temp}
\end{figure*}

% =====================================================================
\section{Inverse Problem}\label{sec:inverse}
% =====================================================================

\subsection{Statement}

Given an observed RGB triple $\bm{R}^{\mathrm{obs}}=(R_{R},R_{G},R_{B})$
and a measured AC modulation amplitude $A^{\mathrm{ac}}_{G}$ from the
green channel relative to a $30\,\mathrm{s}$ rolling baseline, recover
the optical state $\theta=(\melan,\Hbconc,\Soxy,\vaso)$.

\subsection{Identifiability and Channel Separation}

Inspection of \cref{eq:forward} together with
\cref{tab:abs-coeffs} reveals that the four state parameters affect
the three RGB bands with distinct sensitivities:
\begin{itemize}[leftmargin=*,nosep]
\item Melanin dominates the blue band: $\bar{\mu}_{\mathrm{mel},B} L_{2}\!\approx\!0.84$ versus
  $\bar{\mu}_{\Hb,B}\, L_{3,0}\!\approx\!0.084$ per g/dL.
\item Hemoglobin dominates the red band: melanin contribution is small
  ($0.10$) and oxygenation is nearly indistinguishable in red
  ($\bar{\mu}_{\Hb,R}=\bar{\mu}_{\HbO,R}$).
\item Oxygenation acts only in the green band: the difference
  $\bar{\mu}_{\HbO,G}-\bar{\mu}_{\Hb,G}\approx 0.036$ per~g/dL is the only
  observable consequence of $\Soxy$ in the DC reflectance.
\item Vasodilation is captured separately by the AC amplitude.
\end{itemize}
This near-orthogonal sensitivity pattern enables algebraic inversion
without iterative optimisation.

\subsection{Algebraic Inverse}

Define the lifted observations
\begin{equation}
    \ell_{c}
    \;=\;
    \frac{1}{D}\,\frac{R_{c}^{\mathrm{obs}} - \rho_{\mathrm{spec}}}{1 - \rho_{\mathrm{spec}}},
    \qquad c\in\{R,G,B\},
    \label{eq:lift}
\end{equation}
which is the dimensionless product of squared transmittances
$\rho_{2,c}^{2}\rho_{3,c}^{2}$ with the diffuse-return factor removed.
The inversion proceeds in three steps.

\paragraph{Step 1: Melanin from blue.} Approximating the Hb
contribution in blue as a residual constant (justified by the
$\sim$10:1 dominance of melanin in blue), we solve~\cref{eq:rho2} for
$\melan$:
\begin{equation}
    \melan
    \;=\;
    -\frac{\ln \ell_{B}}{2\,\bar{\mu}_{\mathrm{mel},B}\,L_{2}}.
    \label{eq:invert-melanin}
\end{equation}

\paragraph{Step 2: Hemoglobin from red, given melanin.} Correcting the
red lift for the melanin transmittance,
\begin{equation}
    \tilde{\ell}_{R}
    \;=\; \ell_{R}\,\exp(2\,\bar{\mu}_{\mathrm{mel},R}\, \melan\, L_{2}),
    \label{eq:tilde-ellR}
\end{equation}
and noting that $\bar{\mu}_{\Hb,R}=\bar{\mu}_{\HbO,R}$ within band-averaging error,
we solve~\cref{eq:rho3} for $\Hbconc$ at $\vaso=\vaso^{\mathrm{ac}}$
(the value coming from step~4):
\begin{equation}
    \Hbconc
    \;=\;
    -\frac{\ln \tilde{\ell}_{R}}{2\,\bar{\mu}_{\Hb,R}\,L_{3,0}\,\vaso^{\mathrm{ac}}}.
    \label{eq:invert-hb}
\end{equation}

\paragraph{Step 3: Oxygenation from green, given melanin and Hb.} The
green-channel dermal transmittance depends on $\Soxy$ via
\cref{eq:mu-derm}. The corrected green lift
\begin{equation}
    \tilde{\ell}_{G}
    \;=\; \ell_{G}\,\exp(2\,\bar{\mu}_{\mathrm{mel},G}\, \melan\, L_{2})
\end{equation}
satisfies
\begin{equation}
    \tilde{\ell}_{G} \;=\; (1-\Soxy)\, t_{\mathrm{deox}} + \Soxy\, t_{\mathrm{oxy}},
\end{equation}
with
$t_{\mathrm{deox}}=\exp(-2\bar{\mu}_{\Hb,G}\,\Hbconc \vaso^{\mathrm{ac}}\,L_{3,0})$
and analogously for $t_{\mathrm{oxy}}$. Direct inversion gives
\begin{equation}
    \Soxy
    \;=\;
    \frac{\tilde{\ell}_{G} - t_{\mathrm{deox}}}{t_{\mathrm{oxy}} - t_{\mathrm{deox}}}.
    \label{eq:invert-spo2}
\end{equation}

\paragraph{Step 4: Vasodilation from AC amplitude.} The AC modulation
amplitude in the green channel is approximately proportional to the
effective dermal blood volume per cycle,
$A^{\mathrm{ac}}_{G}\propto \Hbconc\,\vaso\,L_{3,0}$. Normalising
against a rolling baseline $A^{\mathrm{base}}_{G}$ and applying a
square-root compression to attenuate noise,
\begin{equation}
    \vaso^{\mathrm{ac}}
    \;=\;
    \sqrt{\max(0.1,\, A^{\mathrm{ac}}_{G}/A^{\mathrm{base}}_{G})}.
    \label{eq:invert-vaso}
\end{equation}
A first-order low-pass with $\tau=60\,\mathrm{s}$ is applied to
$\vaso^{\mathrm{ac}}$ to suppress beat-to-beat noise while preserving
thermoregulatory transients (which evolve on tens of seconds to
minutes).

\begin{figure*}[!htbp]
    \centering
    \includegraphics[width=\textwidth]{figures/panel2_self_consistency.png}
    \caption{\textbf{Self-Consistency of the Two-Iteration Algebraic Inverse Across $1{,}000$ Random Physiological States --- Headline Validation.}
    For each state, the noise-free RGB triple is computed from \cref{eq:forward} and passed back through the two-iteration inverter of \cref{sec:inverse} with $\eta$ supplied as the AC input. Sub-panels A, B, D are 2D scatter plots of recovered against true value, with the identity line in black; sub-panel~C renders the bivariate error structure in 3D.
    (\textbf{A})~Melanin index recovery. Points cluster tightly along the identity line; median relative error $0.5\%$, $99$th-percentile $4.1\%$. The small residual at the top of the range reflects the algebraic approximation in the blue channel where the Hb correction (\cref{eq:tilde-ellR}) leaves a sub-percent residual at high $\melan$.
    (\textbf{B})~Total hemoglobin concentration recovery. Median error $10.1\%$, with a long upper tail at extreme $(\melan, \Hbconc)$ combinations where the band-averaged $\bar{\mu}_{\Hb,R} \approx \bar{\mu}_{\HbO,R}$ approximation breaks down.
    (\textbf{C})~Three-dimensional view of where Hb error concentrates: true melanin (x-axis), true [Hb] (y-axis), and Hb relative recovery error in percent (z-axis, also colour-mapped on the viridis scale). Error grows along the high-melanin wall, identifying the pigment-saturated regime as the recovery's hardest corner.
    (\textbf{D})~Oxygenation recovery: points scatter widely with median $41.2\%$ relative error and visible saturation at the lower clip boundary $\Soxy=0.5$. The behaviour reflects the small green-channel sensitivity to $\Soxy$ visible in \cref{fig:forward-sweeps} and confirms that absolute oxygenation is not recoverable to clinical fidelity from three-band RGB; only relative trends should be used.}
    \label{fig:self-consistency}
\end{figure*}

\subsection{Joint Solution}

Steps 1--4 are coupled only through $\vaso^{\mathrm{ac}}$ entering the
red and green denominators. We initialise with $\vaso=1$, run steps~1
through 3 to estimate $(\melan,\Hbconc,\Soxy)$, then update
$\vaso^{\mathrm{ac}}$ from the AC measurement. A single iteration is
sufficient because the AC channel is nearly orthogonal to the DC
channels under the parameter ranges of interest. The total
computational cost is approximately $50\,\mathrm{FLOP}$ per face
ROI per inversion call, yielding sub-microsecond performance on
contemporary hardware.

% =====================================================================
\section{From Optical State to Cardiac Decomposition}\label{sec:cardiac}
% =====================================================================

\subsection{Skin Temperature}

The temperature mapping is given by~\cref{eq:vaso-T}, applied to the
smoothed vasodilation factor $\vaso^{\mathrm{ac}}$:
\begin{equation}
    \Tskin = 33 + 4\bigl(\vaso^{\mathrm{ac}} - 1\bigr)\,^{\circ}\mathrm{C},
\end{equation}
clipped to $[27, 37]\,^{\circ}\mathrm{C}$.

\subsection{Q$_{10}$ Metabolic Decomposition}

Sino-atrial node automaticity scales with body temperature according
to the Arrhenius $\Qten$~law~\cite{stevens1989q10,delgado2008}:
\begin{equation}
    \mathrm{rate}(T) \;=\; \mathrm{rate}(T_{\mathrm{ref}})\,\Qten^{(T-T_{\mathrm{ref}})/10},
\end{equation}
with $\Qten\approx 2.3$ for cardiac pacemaker activity and
$T_{\mathrm{ref}}=37\,^{\circ}\mathrm{C}$. Linearising about the
resting forehead reference $T_{\mathrm{rest}}=33\,^{\circ}\mathrm{C}$,
the metabolic shift from baseline is
\begin{equation}
    \dHRm \;=\; \alpha_{T}\,(\Tskin - T_{\mathrm{rest}})\,\HRint,
    \quad
    \alpha_{T} \approx 0.08\,\mathrm{K}^{-1},
    \label{eq:dHR-met}
\end{equation}
where $\alpha_{T}$ derives from $\Qten=2.3$. The hypoxic compensation
acts via chemoreceptor reflexes:
\begin{equation}
    \dHRo \;=\; \beta_{\mathrm{O_2}}\,(1 - \Soxy)\,\HRint,
    \quad
    \beta_{\mathrm{O_2}}\approx 0.15,
    \label{eq:dHR-o2}
\end{equation}
with $\beta_{\mathrm{O_2}}$ calibrated from chemoreceptor-driven
heart-rate elevation per $10\%$ drop in
$\SpO$~\cite{guyenet2014}. The autonomic residual is then
\begin{equation}
    \boxed{\;
    \dHRa \;=\; \HRobs - \HRint - \dHRm - \dHRo,
    \;}
    \label{eq:dHR-auto}
\end{equation}
which isolates the genuine neural drive (sympathovagal balance) from
purely thermal and hypoxic contributions. The intrinsic rate
$\HRint$ is set to a baseline of $60\,\mathrm{bpm}$ in the absence of
calibration data; subject-specific calibration via the lowest sustained
HR over a 7-day window is straightforward.

\subsection{Effect on Inotropy Estimation}

A critical consequence of the decomposition is that the autonomic
component, not the raw HR, should drive any inotropy estimator. In
ventricular elastance models the end-systolic elastance $E_{\mathrm{es}}$
scales with sympathetic outflow, which is well represented by
$\dHRa$ but not by $\HRobs$ when the subject is warm or
hypoxic~\cite{sagawa1988,burkhoff2007}. Substituting $\dHRa$ for
$\HRobs$ in any standard ventricular-elastance update rule
straightforwardly removes thermal-drive contamination from
contractility estimates.

% =====================================================================
\section{Implementation}\label{sec:implementation}
% =====================================================================

The full pipeline runs in real time in a WebGPU-enabled browser. A
720p front-facing camera frame is processed per render frame ($30$~Hz
typical) as follows.

\paragraph{Face localisation.} A lightweight face landmarker
(\textsc{MediaPipe} \texttt{FaceLandmarker} v0.10
\cite{mediapipe2019}) extracts 478 face landmarks per frame in
$\sim$5\,ms on commodity GPU. Forehead, left-cheek, and right-cheek
ROIs are constructed as inset bounding boxes around stable index
subsets.

\paragraph{Per-pixel rPPG.} A WebGPU compute shader downsamples the
forehead ROI to a $64\times 64$ patch grid, computes per-patch RGB
mean, and appends the green channel into a circular history buffer of
$256$ frames ($\sim$8.5\,s). A second compute pass computes a
detrended autocorrelation in the cardiac band (0.7--3\,Hz) per patch,
yielding per-pixel $R_{c}$ coherence, AC amplitude, period, and SNR
estimates as a $64\times 64$ \texttt{rgba16float} field. The face
ROI is sampled separately on the CPU via a small offscreen 2D canvas
to obtain mean RGB at $\sim$5\,Hz.

\paragraph{Inversion.} The algebraic inverter
(\cref{eq:invert-melanin,eq:tilde-ellR,eq:invert-hb,eq:invert-spo2,eq:invert-vaso})
runs once per second on the most recent face-ROI mean RGB and the AC
amplitude, producing the optical state $\theta(t)$ and the derived
quantities $(\Tskin, \dHRm, \dHRo, \dHRa)$. State estimates are
exponentially smoothed with $\alpha=0.18$ to suppress per-second
noise without lagging thermoregulatory transients.

\paragraph{Persistence and dashboard.} Per-second records are
auto-saved to an in-browser \textsc{IndexedDB} store and exposed
through a brushable \textsc{D3} + \textsc{Crossfilter}~\cite{bostock2011d3}
dashboard, supporting longitudinal sessions of $\geq 7$ days
without back-end infrastructure.

\paragraph{Computational budget.} The dominant cost is MediaPipe
landmark inference ($\sim$5\,ms/frame) and the WebGPU shader pipeline
($\sim$1\,ms/frame). The inversion itself is sub-microsecond per call,
running at $1\,\mathrm{Hz}$. Total CPU utilisation on a typical
laptop is $\sim$8\%.



% =====================================================================
\section{Experimental Methods}\label{sec:methods}
% =====================================================================

We validate the method through six simulation experiments. Direct
in vivo validation against laboratory-grade reference instruments is
beyond the scope of this paper but is the natural next step
(\cref{sec:discussion}). All experiments are deterministic and
reproducible from the supplied source code; results are saved as JSON
and rendered to PNG panels.

\subsection{Experiment 1: Forward-Model Spectra}

We tabulate the predicted RGB reflectance as a function of each state
parameter individually, holding the others at the reference values
$\theta_{0}=(\melan{=}0.15,\, \Hbconc{=}14\,\mathrm{g/dL},\,
\Soxy{=}0.97,\, \vaso{=}1.0)$. Each parameter is swept across its
physiological range with $128$ samples; the forward model
(\cref{eq:forward}) is evaluated and the three channel responses
plotted. This experiment characterises the sensitivity of each
observable channel to each state parameter.

\begin{figure*}[!htbp]
    \centering
    \includegraphics[width=\textwidth]{figures/panel1_forward_sweeps.png}
    \caption{\textbf{Forward-Model Sensitivity per State Parameter and the Joint RGB Manifold --- Channel-Separation Diagnostic.}
    All sub-panels evaluate the layered Beer--Lambert forward map of \cref{eq:forward} with channel-band-averaged absorption coefficients from \cref{tab:abs-coeffs}. In sub-panels~A--C the three remaining state parameters are held at the reference $\theta_{0}=(\melan{=}0.15,\, \Hbconc{=}14\,\mathrm{g/dL},\, \Soxy{=}0.97,\, \vaso{=}1.0)$.
    (\textbf{A})~Reflectance versus melanin index over $[0.02, 0.85]$. Blue collapses from $\sim 0.09$ to $\sim 0.05$, green from $\sim 0.05$ to $\sim 0.04$, and red from $\sim 0.42$ to $\sim 0.36$. The blue-channel sensitivity to melanin is $\sim 8\times$ the red-channel sensitivity after the round-trip squaring; this dominance underwrites step~1 of the inverter (\cref{eq:invert-melanin}).
    (\textbf{B})~Reflectance versus total hemoglobin concentration over $[8, 20]\,\mathrm{g/dL}$. Red drops $\sim 0.46\!\to\!0.36$, green $\sim 0.08\!\to\!0.04$, blue $\sim 0.15\!\to\!0.05$. The red-channel response is the cleanest because melanin contribution is small there and oxygenation has zero effect at band-averaged level, supporting step~2 of the inverter (\cref{eq:invert-hb}).
    (\textbf{C})~Reflectance versus vasodilation factor $\eta$ over $[0.6, 1.6]$. All channels respond because $\eta$ amplifies the dermal-blood optical path; the response is most prominent in green and is what makes the AC-amplitude readout in \cref{eq:invert-vaso} sensitive to thermoregulatory transients.
    (\textbf{D})~Three-dimensional joint cloud of $(R, G, B)$ reflectance for $1{,}500$ states sampled uniformly across the physiological ranges, points coloured by oxygenation $\Soxy$ on a plasma colour map. The cloud occupies a low-dimensional curved manifold in RGB-space; oxygenation varies along the small green-axis component visible in the colour gradient and is the single dimension along which the inverse is most poorly conditioned.}
    \label{fig:forward-sweeps}
\end{figure*}

\subsection{Experiment 2: Self-Consistency}

We sample $N=1{,}000$ random states $\theta_{i}$ uniformly within the
physiological ranges
$\melan\in[0.05,0.85]$, $\Hbconc\in[8,20]\,\mathrm{g/dL}$,
$\Soxy\in[0.85,1.0]$, $\vaso\in[0.6,1.6]$. For each $\theta_{i}$ we
compute the noise-free RGB
$\bm{R}_{i}=F(\theta_{i})$ and apply the inversion to recover an
estimate $\hat{\theta}_{i}$, supplying the true $\vaso_{i}$ as the AC
input. The relative error per parameter is reported as
$|\hat{\theta}_{i,k}-\theta_{i,k}|/\theta_{i,k}$. Median, $90$th, and
$99$th percentile errors quantify the algebraic-inverse fidelity in
the absence of noise.

\subsection{Experiment 3: Noise Sensitivity}

For each SNR value
$\mathrm{SNR}\in\{10,15,20,25,30,35,40,45,50\}\,\mathrm{dB}$, we
generate $N=200$ random states $\theta_{i}$, compute noise-free RGB,
and add zero-mean Gaussian noise to each channel scaled to the target
SNR relative to the channel's signal power. We invert the noisy RGB
and compute the RMS relative error per parameter. The resulting
curves characterise graceful degradation under measurement noise.

\subsection{Experiment 4: Vasodilation--Temperature Mapping}

We evaluate \cref{eq:vaso-T} across the vasodilation range and overlay
empirically observed forehead skin-temperature ranges from the
thermoregulation literature: cold-pressor minimum, quiet-sitting
median, mild-exercise plateau, and maximum-exercise upper bound. The
panel illustrates the physiological consistency of the linearised
mapping.

\subsection{Experiment 5: PCHR Scenario Attribution}

We define four synthetic scenarios with known dominant drives:
\begin{enumerate}[label=(\alph*),nosep]
\item \textbf{Rest:} $\Tskin=33\,^{\circ}\mathrm{C}$,
  $\Soxy=0.97$, $\HRobs=70$.
\item \textbf{Fever:} $\Tskin=35.5\,^{\circ}\mathrm{C}$,
  $\Soxy=0.97$, $\HRobs=85$.
\item \textbf{Hypoxia:} $\Tskin=33\,^{\circ}\mathrm{C}$,
  $\Soxy=0.85$, $\HRobs=85$.
\item \textbf{Exercise:} $\Tskin=34.5\,^{\circ}\mathrm{C}$,
  $\Soxy=0.96$, $\HRobs=120$.
\end{enumerate}
For each scenario we apply
\cref{eq:dHR-met,eq:dHR-o2,eq:dHR-auto} with $\HRint=60\,\mathrm{bpm}$
and report the four components. The expected pattern is dominance of
$\dHRm$ in (b), $\dHRo$ in (c), and $\dHRa$ in (d), with negligible
contributions of all non-baseline drives in (a).

\begin{figure*}[!htbp]
    \centering
    \includegraphics[width=\textwidth]{figures/panel5_pchr_scenarios.png}
    \caption{\textbf{PCHR Decomposition: Scenario Attribution and the Underlying Decomposition Surfaces --- Headline Diagnostic Output.}
    Each sub-panel exercises the four-component decomposition $\HRobs = \HRint + \dHRm + \dHRo + \dHRa$ from \cref{eq:dHR-met,eq:dHR-o2,eq:dHR-auto} with $\HRint = 60\,\mathrm{bpm}$ and the normoxic reference $\Soxy_{\mathrm{ref}} = 0.97$.
    (\textbf{A})~Stacked-bar decomposition for the four canonical scenarios (rest, fever, hypoxia, exercise) defined in \cref{tab:pchr-scenarios}. Bottom segments (grey) are the intrinsic baseline of $60\,\mathrm{bpm}$; orange the metabolic drive $\dHRm$; green the hypoxic compensation $\dHRo$; purple the autonomic residual $\dHRa$. Each scenario's primary drive is correctly identified by the dominant non-baseline contribution: thermal in fever ($\dHRm=12\,\mathrm{bpm}$), hypoxic in hypoxia ($\dHRo=10.8\,\mathrm{bpm}$), autonomic in exercise ($\dHRa=51.9\,\mathrm{bpm}$).
    (\textbf{B})~Metabolic drive $\dHRm$ as a continuous function of $\Tskin$ at $\HRint = 60\,\mathrm{bpm}$, with the four scenario points (\textit{rst}, \textit{fev}, \textit{hyp}, \textit{exe}) overlaid on the curve. Fever and exercise sit above the rest line at the temperatures their scenarios specify; hypoxia sits exactly at zero because its $\Tskin = 33\,^{\circ}\mathrm{C}$ is the resting reference.
    (\textbf{C})~Three-dimensional decomposition surface: $\dHRa = \HRobs - \HRint - \dHRm(\Tskin) - \dHRo(\Soxy)$ across the $(\Tskin, \Soxy)$ plane at a representative $\HRobs = 100\,\mathrm{bpm}$ on a coolwarm colour map. Higher temperature and lower oxygenation both reduce the autonomic residual because both terms consume part of the observed-HR budget. The four scenario points (black markers) sit at heights corresponding to their specific $\HRobs$ values.
    (\textbf{D})~Hypoxic drive $\dHRo$ as a continuous function of $\Soxy$, with the four scenario points overlaid. The drive is exactly zero at and above the normoxic reference $\Soxy = 0.97$; the hypoxia scenario at $\Soxy = 0.85$ produces $10.8\,\mathrm{bpm}$ of chemoreflex-driven HR elevation.}
    \label{fig:pchr-scenarios}
\end{figure*}

\subsection{Experiment 6: Comparison with Conventional Green-Only rPPG}

To quantify the lighting-robustness benefit of the layered inversion
relative to amplitude-only rPPG, we simulate a $5$-minute recording
under a synthetic time-varying ambient illumination drift of
$\pm 20\%$ (sinusoidal at $0.05\,\mathrm{Hz}$). Ground-truth state is
held constant at $\theta_{0}$ with $\HRobs=70$, and the green-channel
DC is multiplied by the drift envelope. We compare two heart-rate
estimators:
\begin{enumerate}[label=(\alph*),nosep]
\item \textbf{Green-only baseline:} bandpass-filtered green-channel
  amplitude, peak-detected, converted to RR and HR.
\item \textbf{Layered inversion:} the same green-channel amplitude
  fed through the inversion to recover $\vaso$ separately from the DC
  drift, then HR derived from the AC peaks of the recovered $\vaso$
  time series.
\end{enumerate}
The variance of the resulting HR estimates over the $5$-minute window
is the comparison metric.

\subsection{Reproducibility}

All experiments are implemented in Python~3 with NumPy~1.26 and
Matplotlib~3.8. The forward model is implemented to numerical
equivalence with the TypeScript implementation used in the real-time
browser tool. JSON outputs and PNG panels are deterministic given
fixed random seeds. Full source is available alongside this paper.

% =====================================================================
\section{Results}\label{sec:results}
% =====================================================================

\subsection{Forward-Model Spectra (Experiment 1)}

\Cref{fig:forward-spectra} shows the predicted RGB reflectance as
each of the four state parameters is swept individually. The
sensitivity pattern matches the channel-separation strategy of the
inverse problem: the blue channel responds primarily to melanin, the
red and green channels respond to hemoglobin and vasodilation, and
oxygenation acts visibly only on green. The near-orthogonality of
sensitivity is what makes the algebraic inverse well-conditioned.



\subsection{Self-Consistency (Experiment 2)}

\Cref{tab:self-consistency} reports the recovery error percentiles
for each parameter under noise-free conditions, computed over
$1{,}000$ random physiological states sampled uniformly from the
ranges $\melan\in[0.05,0.85]$, $\Hbconc\in[8,20]\,\mathrm{g/dL}$,
$\Soxy\in[0.85,1.0]$, $\vaso\in[0.6,1.6]$. The two-iteration
inverse converges to median errors of $0.5\%$ for melanin and
$10\%$ for total hemoglobin, with the $90$th and $99$th
percentiles indicating how the algebraic approximations degrade at
the extremes of the parameter ranges. Oxygenation
recovery is structurally limited: the green-channel difference
between oxy- and deoxyhemoglobin
($\bar{\mu}_{\HbO,G}-\bar{\mu}_{\Hb,G}\approx 0.036$ per~g/dL of
Hb) is small relative to the total green-channel attenuation, and
the inversion is therefore highly sensitive to upstream estimation
of melanin and Hb. The resulting median oxygenation error of
$\approx 41\%$ is essentially independent of measurement noise (see
\cref{ssec:noise}) and reflects an intrinsic limitation of three-band
RGB spectroscopy for oxygenation. We discuss this further in
\cref{sec:discussion}.

\begin{table}[H]
\centering
\caption{Self-consistency: relative recovery error percentiles over
$1{,}000$ random physiological states (noise-free, two iterations of
the algebraic inverse).}
\label{tab:self-consistency}
\small
\begin{tabular}{lccc}
\toprule
Parameter & Median & 90th \% & 99th \% \\
\midrule
Melanin index $\melan$ & 0.5\% & 1.0\% & 4.1\% \\
Total Hb $\Hbconc$ & 10.1\% & 26.8\% & 47.1\% \\
Oxygenation $\Soxy$ & 41.2\% & 48.3\% & 49.8\% \\
Vasodilation $\vaso$ & 0\% (exact) & 0\% & 0\% \\
\bottomrule
\end{tabular}
\end{table}



\subsection{Noise Sensitivity (Experiment 3)}\label{ssec:noise}

\Cref{fig:noise-sensitivity} plots the RMS recovery error per
parameter against the additive-noise SNR. Two qualitatively
different regimes are visible. Melanin and total hemoglobin are
\emph{noise-limited}: their RMS errors fall as SNR rises, reaching
roughly $1\%$ and $16\%$ respectively at $50\,\mathrm{dB}$.
Oxygenation is \emph{model-limited}: the RMS error stays
approximately flat near $35\%$ across all SNR values, confirming
that the floor on oxygenation recovery is set by the band-averaged
spectral approximation in \cref{eq:mu-derm}, not by measurement
noise. Improving oxygenation accuracy therefore requires
spectrally-narrower data than RGB provides (e.g., a fourth band
isolating the $660$/$940\,\mathrm{nm}$ oxy/deoxy crossover) rather
than better noise rejection.

\begin{figure*}[!htbp]
    \centering
    \includegraphics[width=\textwidth]{figures/panel3_noise_sensitivity.png}
    \caption{\textbf{Noise Sensitivity: Two Distinct Regimes (Noise-Limited vs.\ Model-Limited) Across the Three Recovered Parameters.}
    Each curve is computed over $200$ random physiological states with channel-wise zero-mean Gaussian noise scaled to the target SNR; the summary statistic is RMS relative recovery error per parameter. The dashed reference line at $5\%$ marks a practical usability threshold.
    (\textbf{A})~Melanin RMS error versus SNR (log y-axis). The error decreases gradually with rising SNR but stays above the $5\%$ usability line across the entire range tested. Recovery is dominated by noise-induced clipping events at the boundaries of the parameter range when noise pushes the blue-channel reflectance outside its physiologically plausible support and the algebraic step~1 saturates.
    (\textbf{B})~Hemoglobin RMS error versus SNR (log y-axis). Smooth monotonic decline from $\sim 50\%$ at $10\,\mathrm{dB}$ to $\sim 16\%$ at $50\,\mathrm{dB}$. This is the canonical \emph{noise-limited} regime: the recovery improves directly with measurement-noise reduction and would benefit further from longer-window averaging or a low-noise camera.
    (\textbf{C})~Three-dimensional surface over (SNR, parameter index, RMS error in percent). The surface separates the three parameters into visibly different regimes: a sloped melanin slice rising at low SNR, a smooth Hb descent across the SNR range, and a flat SpO$_{2}$ slice that the colour mapping picks out as a model-limited plateau.
    (\textbf{D})~SpO$_{2}$ RMS error versus SNR (linear y-axis). Approximately constant at $35$--$40\%$ across all SNR values --- the canonical \emph{model-limited} regime. The error floor is set by the band-averaged spectral approximation in \cref{eq:mu-derm}, not by measurement noise. No noise-side improvement would change this floor; reducing it requires spectrally narrower data, e.g.~a fourth band isolating the $\sim 660/940\,\mathrm{nm}$ oxy/deoxy crossover that conventional pulse oximetry exploits.}
    \label{fig:noise-sensitivity}
\end{figure*}



\subsection{Vasodilation--Temperature Mapping (Experiment 4)}

\Cref{fig:vaso-temp} shows the linear mapping and overlays the
empirically observed temperature ranges. The mapping spans the
realistic $[27, 37]\,^{\circ}\mathrm{C}$ range over the
vasodilation interval $[0.6, 1.6]$, with the resting state
($\vaso=1$) corresponding to $33\,^{\circ}\mathrm{C}$ as
intended.



\subsection{PCHR Scenario Attribution (Experiment 5)}

\Cref{fig:pchr-scenarios} shows the decomposition of $\HRobs$ into
the four components for each of the four scenarios.
\Cref{tab:pchr-scenarios} reports the numerical breakdown. In the
rest scenario all non-intrinsic components are zero except the
small autonomic residual ($10\,\mathrm{bpm}$) consistent with
typical resting parasympathetic withdrawal. In fever, the metabolic
component contributes $12\,\mathrm{bpm}$, accounting for the
$\Qten$-driven thermal acceleration of the sino-atrial pacemaker;
the autonomic residual remains modest. In hypoxia, the hypoxic
component contributes $10.8\,\mathrm{bpm}$ at $\Soxy=0.85$ via the
chemoreceptor reflex; the autonomic residual at $14.2\,\mathrm{bpm}$
captures the additional sympathetic outflow that hypoxia recruits
beyond the linear chemoreflex term. In exercise, the autonomic
residual ($51.9\,\mathrm{bpm}$) dominates as expected, with a
secondary metabolic contribution ($7.2\,\mathrm{bpm}$) from the mild
thermal elevation. The decomposition correctly attributes each
scenario's primary drive to the matching component.



\begin{table}[H]
\centering
\caption{PCHR decomposition results for the four synthetic
scenarios. All values in bpm. Intrinsic baseline is $60\,\mathrm{bpm}$
in all scenarios; the SpO$_{2}$ reference for the hypoxic
component is $0.97$, so normoxic ($\Soxy=0.97$) gives
$\dHRo=0$.}
\label{tab:pchr-scenarios}
\small
\begin{tabular}{lcccc}
\toprule
Scenario & $\HRobs$ & $\dHRm$ & $\dHRo$ & $\dHRa$ \\
\midrule
Rest     & 70  &  0.0 &  0.0 & 10.0 \\
Fever    & 85  & 12.0 &  0.0 & 13.0 \\
Hypoxia  & 85  &  0.0 & 10.8 & 14.2 \\
Exercise & 120 &  7.2 &  0.9 & 51.9 \\
\bottomrule
\end{tabular}
\end{table}

\subsection{Comparison with Conventional rPPG (Experiment 6)}

\Cref{fig:conventional-comp} (left) plots HR estimation variance
against the amplitude of slow ($0.05\,\mathrm{Hz}$) ambient-light
drift for both methods. Both perform comparably across the swept
range: the cardiac-band band-pass cleanly rejects the slow drift in
both cases, and the layered method's additional DC-normalisation
step contributes a small extra noise term that slightly raises its
variance at high drift levels. HR alone, therefore, does not
distinguish the methods under typical lighting drift conditions.

The right panel reports the metric where the layered method does
unambiguously outperform: SpO$_{2}$ recovery error from the same
camera under the same drift. The conventional green-only baseline
produces no SpO$_{2}$ estimate at all; the layered method's full
state recovery yields SpO$_{2}$ within $\approx 3\%$ relative error
of the true value across all drift amplitudes, because the
DC-normalisation step removes the drift envelope from the RGB triple
before inversion. The same DC-normalisation pipeline that adds
modest noise to HR variance is what makes additional state recovery
robust under drift.

The honest summary is that the layered method's distinct value is
\emph{not} a smaller HR variance under typical conditions but the
recovery of state parameters that conventional rPPG does not produce
at all: vasodilation, skin temperature, total hemoglobin, and a
relative oxygenation trend. These additional signals are what feed
the downstream PCHR decomposition (\cref{sec:cardiac}) and the
metabolic-autonomic separation that gives the layered approach
qualitatively different diagnostic content.

\begin{figure*}[!htbp]
    \centering
    \includegraphics[width=\textwidth]{figures/panel6_conventional_comparison.png}
    \caption{\textbf{Comparison with Conventional Green-Only rPPG Across a Drift-Amplitude Sweep --- Honest HR Parity Plus Unique State Recovery.}
    A 5-minute synthetic recording is generated at each of seven slow-drift amplitudes (0--100\% of DC, $0.05\,\mathrm{Hz}$ ambient envelope) on a baseline cardiac AC modulation of $0.6\%$ peak-to-peak in green. Both methods share an identical Hann-windowed FFT peak-picker in the $0.7$--$3.0\,\mathrm{Hz}$ cardiac band; the only difference is the front-end signal (raw $G$ for method~A vs.~the DC-normalised ratio $G/G_{\mathrm{DC}}$ for method~B).
    (\textbf{A})~Heart-rate estimation variance versus drift amplitude. At low drift ($\leq 20\%$) both methods perform near-identically. As drift grows, method~B's DC-normalisation step introduces a small additional noise term and its variance rises faster than method~A's; this is the honest finding that under bandpass-rejectable drift the layered method is not an HR-variance winner.
    (\textbf{B})~Heart-rate bias versus drift amplitude. Both methods exhibit small positive bias growing with drift, with method~B tracking method~A closely up to $\sim 40\%$ drift and then diverging upward as residual fluctuations from the DC-normalisation step couple into the spectral-peak picker.
    (\textbf{C})~Three-dimensional view of method~B's HR-estimate manifold across (drift amplitude, time, HR estimate) on a viridis colour map. Each coloured trace is the per-window HR estimate at one drift level; lower-drift traces (purple) sit cleanly at the true 70-bpm plane while higher-drift traces (yellow) show short-window excursions at the FFT-bin resolution. The 3D rendering exposes the temporal structure that the variance summary in panel~A flattens out.
    (\textbf{D})~SpO$_{2}$ recovery error of method~B's full state inversion versus drift amplitude. The error is essentially constant near the structural-floor value identified in \cref{fig:noise-sensitivity} regardless of drift amplitude, because the time-averaged DC normalisation removes the drift envelope from the RGB triple before inversion. The conventional green-only baseline produces no SpO$_{2}$ estimate at any drift level. This panel anchors the honest summary of the entire experiment: the layered method's distinct value is not improved HR variance but the existence of additional channels (vasodilation, $\Tskin$, $[\Hb]$, relative SpO$_{2}$ trend) at all.}
    \label{fig:conventional-comparison}
\end{figure*}


% =====================================================================
\section{Discussion}\label{sec:discussion}
% =====================================================================

\subsection{What the Method Does and Does Not Provide}

The method transforms an rPPG pipeline from amplitude extraction into
a model-grounded biophysical inversion. Concretely it provides:
(i) a melanin index that varies across subjects and is invariant to
ambient light; (ii) a total hemoglobin concentration estimate;
(iii) an oxygenation proxy; (iv) a vasodilation factor from which a
skin-temperature estimate follows; and (v) a four-way decomposition
of heart rate into intrinsic, metabolic, hypoxic, and autonomic
components. It does \emph{not} provide an absolute pulse-oximetry-grade
$\SpO$ reading---the green-channel sensitivity to oxygenation is
small relative to the bandpass-averaging error in our coefficients,
and conventional pulse oximetry uses far more spectrally selective
$660$/$940\,\mathrm{nm}$ pairs not available on a colour camera.
The oxygenation channel is therefore best understood as a model
proxy whose relative trends are informative but whose absolute
calibration requires a separate reference.

\subsection{Limitations}

\paragraph{Coefficient calibration.} The band-averaged absorption
coefficients in \cref{tab:abs-coeffs} are computed from population
spectra and are not subject-specific. Individual variation in skin
chromophore concentrations, capillary geometry, and surface
microstructure introduces systematic bias. Subject-specific
calibration via the lowest sustained skin temperature over a 7-day
window (an intrinsic forehead $\Tskin$) is straightforward and
recommended for clinical applications.

\paragraph{Motion artefacts.} The forward model assumes a stationary
face ROI under approximately constant illumination geometry. Head
motion causes ROI mis-registration and changes the specular
geometry, both of which produce out-of-model variation that the
inverter will absorb into spurious state changes.
Motion-compensation pre-filters (e.g.~CHROM-style chrominance
projection) can be applied upstream of the inversion if needed.

\paragraph{Extreme pigmentation.} For very dark skin
($\melan>0.7$), the blue-channel signal is heavily attenuated by
melanin and the SNR for the melanin estimate itself is
reduced. The algebraic inverter degrades gracefully but the
oxygenation estimate is more sensitive to noise in this regime.

\paragraph{Temperature absolute calibration.} The
vasodilation--temperature linearisation (\cref{eq:vaso-T}) is
calibrated from the published thermoregulation literature, not from
per-subject reference thermometry. Absolute $\Tskin$ readings
should be treated as model proxies; relative changes (warming,
cooling transients) are robust.

\subsection{Future Directions}

\paragraph{In vivo validation.} The natural next step is direct
validation against laboratory reference instruments: ECG for HR
ground truth; pulse oximetry for $\SpO$; infrared thermography for
$\Tskin$; and cuff oscillometry or arterial line for blood pressure.
A protocol involving graded thermal challenges (cold pressor, warm
bath) and ventilatory challenges (breath-hold, hypoxic gas) will
exercise the metabolic and hypoxic components independently and
provide quantitative bounds on the method's clinical utility.

\paragraph{Active illumination.} The screen itself is a
spectrally configurable light source whose output is exactly known.
Brief modulated calibration sequences (e.g., $1\,\mathrm{Hz}$ red /
green / blue alternation at low amplitude) can be projected onto the
user's face during inactive moments and the reflected response
inverted to obtain channel-specific reflectance gains, removing the
need for population-average constants in
\cref{tab:abs-coeffs}.

\paragraph{Multi-region cross-validation.} The forehead, both
cheeks, and the nose all carry the same cardiac modulation but with
slightly different layer thicknesses and capillary densities.
Joint inversion across all four ROIs constrains the inverse problem
beyond what any single ROI provides and exposes regional
inconsistencies as a model-validity diagnostic.

\paragraph{Cross-channel consistency residual.} The forward model
predicts a particular relationship between R, G, and B for any
admissible state. The residual
$\|\bm{R}^{\mathrm{obs}}-F(\hat{\theta})\|$ is a
real-time scalar that flags when the model is locally invalid (e.g.
under heavy motion or unusual lighting). Surfacing this residual to
the user allows automatic confidence estimation per measurement.

% =====================================================================
\section{Conclusion}\label{sec:conclusion}
% =====================================================================

We have presented a forward-model-based inversion of camera-measured
forehead skin reflectance that recovers melanin, hemoglobin
concentration, oxygenation, and a vasodilation factor from RGB pixel
data and the AC modulation of the green channel. The inverse is
algebraic, runs in sub-microsecond time per call, and replaces the
pixel-blind amplitude extraction of conventional rPPG with a
biophysically grounded multi-parameter estimate. From the recovered
state we derive a skin-temperature estimate via thermoregulatory
linearisation and a four-way decomposition of heart rate into
intrinsic, metabolic, hypoxic, and autonomic components. Simulation
validation across self-consistency, noise sensitivity, temperature
mapping, scenario attribution, and a comparison with conventional
green-only rPPG confirms the method's correctness within the
inherent limits of three-band RGB spectroscopy. Melanin and total
hemoglobin recover with median errors of $0.5\%$ and $10\%$
respectively; oxygenation has a structural ceiling near $40\%$
relative error and is best treated as a relative-trend signal.
Under bandpass-rejectable ambient-light drift the layered method's
heart-rate variance is comparable to a green-only baseline; the
distinct value of the method is the recovery of vasodilation, skin
temperature, total hemoglobin, and a metabolic--autonomic
decomposition of heart rate that amplitude-only rPPG does not
produce. The full pipeline runs in real time in a WebGPU-enabled
browser. The natural next steps are in vivo validation against
reference instrumentation and per-subject calibration via brief
active illumination sequences, both of which are expected to
substantially improve the absolute fidelity of the recovered state
beyond the population-average band-coefficient regime examined here.

\bibliographystyle{unsrt}
\bibliography{references}

\end{document}